\section{总结与优化方向}\label{sec:conclusion}

\subsection{已完成工作}

截至本报告撰写时，系统已完成三层记忆架构（短期 / 中期 / 长期）与双库（偏好 / 知识）
落地，全部状态文件化、零外部服务依赖；实现多源偏好提取三路提取器、只追加版本化、
非破坏回滚、按时间点回溯与基于场景相似度的跨场景复用；建成 SPO 事实版本链、同义
归一化、六动作冲突处理、关联加权检索与证据先行的模板生命周期；落实六类敏感信息检测、
全入库边界脱敏与两阶段精准遗忘（预案—审批—物理删除—残留复验—审计）；交付可插拔
嵌入三后端（remote / local / kylin）、int8 量化向量库与 v1/v2 格式迁移；记忆子系统
已接入 agent 主循环，11 个集成指标测试与数十个模块单元测试通过（2026-09-13 运行）。

\subsection{待完成事项}

以下事项均在前文以红色 TODO 标注，此处按工作类别汇总。（1）\todo{标准化评测数据集
三子集构建与冻结（\S\ref{sec:dataset}）}；（2）\todo{四项性能指标实测与消融实验
（\S\ref{sec:metrics}）}；（3）\todo{检索延迟网格实测（\S\ref{sec:metrics}）}；
（4）\todo{用户行为 / 手动配置管道接入 agent 主循环（\S\ref{sec:pref}）}；
（5）\todo{偏好注入系统提示词，形成个性化闭环演示（\S\ref{sec:pref}）}；
（6）\todo{中期 $\rightarrow$ 长期晋升与中期淘汰机制（\S\ref{sec:flow}）}；
（7）\todo{麒麟 embedding SDK 真机联调（\S\ref{sec:kylin-sdk}）}；（8）\todo{银河
麒麟桌面 OS 适配测试报告（\S\ref{sec:deploy}）}；（9）\todo{真实场景案例实测记录与
演示视频（\S\ref{sec:case}）}；（10）\todo{LoCoMo 数据集按 Recall 口径重测
（\S\ref{sec:evaluation}）}。

\subsection{后续优化方向}

后续优化沿四个方向推进。\textbf{本地神经嵌入}：以小型中文嵌入模型（int8 量化）替换
哈希基线，提升语义召回上限，同时保持离线与低资源占用。\textbf{混合检索排序}：结构化
加权检索与向量近邻的融合排序，兼顾精确字段匹配与语义泛化，知识库规模增长后引入轻量
ANN 索引。\textbf{遗忘意图的 NLU 扩展}：在保留规则式安全护栏的前提下，扩展意图解析
的概念覆盖面。\textbf{中期记忆生命周期}：重要性衰减、访问计数驱动晋升与淘汰，补全
三层流转的最后一环。
